%MAIN-DOCUMENT
%INCLUDIERT ALLER FILES DER ORDNERSTRUKTUR

% Header beinhaltet Dokumentenklasse sowie includierte Packages
%HEADER

%<---!!!!!!!!!!!!!!! MAKRO-DEFIONITIONEN; BITTE NICHT VERAENDERN !!!!!!!!!!
%<--- ARBEIT EINSEITIG
\def\makroEinseitig{
%KOMA-Script-Klasse: scrreprt
%deutsches Design, Schriftgröße 12, DIN A4
%Literaturverzeichnis und Index in Inhaltsverzerzeichnis einbinden
\documentclass[12pt,a4paper,listof=totoc,oneside]{scrreprt}
%Seitenspiegel einstellen
\usepackage[a4paper]{geometry}
\geometry{a4paper,left=30mm,right=25mm,
bottom=20mm,top=15mm,bindingoffset=2mm,
includehead,includefoot}}
% ARBEIT EINSEITIG --->

\def\makroZweiseitig{
%<--- ARBEIT ZWEISEITIG
%KOMA-Script-Klasse: scrreprt
%deutsches Design, zweiseitig
%Literaturverzeichnis und Index in Inhaltsverzerzeichnis einbinden
\documentclass[12pt,a4paper,listof=totoc,twoside, headsepline]{scrreprt}
\usepackage[a4paper]{geometry}
\geometry{a4paper,left=25mm,right=25mm,
bottom=20mm,top=15mm,bindingoffset=2mm,
includehead,includefoot}}
% ARBEIT ZWEISEITIG --->

%<--- Einstellungen Kopfzeile
\def\makroFH-Kopfzeilenstil{
\pagestyle{scrheadings} 
\setheadsepline{0.4pt}
\pagestyle{scrheadings}
\renewcommand*{\chapterpagestyle}{scrheadings}}
%Einstellungen Kopfzeile --->
%!!!!!!!!!!!!!!! MAKRO-DEFIONITIONEN; BITTE NICHT VERAENDERN !!!!!!!!!!--->


%AUSWAHL: TEXT EINSEITIG (ja/nein)
\makroEinseitig
%\makroZweiseitig

%schalte Umlaute frei
\usepackage[ngerman]{babel}
%passende Codierung
\usepackage[utf8]{inputenc}
%Seitenspiegel einzustellen
\usepackage[a4paper]{geometry}
%Mathepaket
\usepackage{amsmath}
%Symbole
\usepackage{amssymb}
%griechische Symbole
\usepackage{upgreek}
%weitere Symbole
\usepackage{pxfonts}
% Phonetischen Alphabete für LaTeX
\usepackage{tipa}
%farbige Schriften
\usepackage{color}
\usepackage{scrhack}
%Bilder fixieren
\usepackage{float}
%Grafiken einbinden
\usepackage{graphicx}
% Kopf- und Fußzeilen
\usepackage[automark,standardstyle,markusedcase]{scrpage2}
% deutsche Überschriften
\usepackage[ngerman]{translator}
% Kopfzeilenabstand festlegen
\setlength{\headheight}{10mm}
%Abb. statt Abbildung
\usepackage{caption3}
\addto\captionsngerman{
\renewcommand{\figurename}{Abb.}
\renewcommand{\tablename}{Tab.}
}
%Glossar-Pakage
\usepackage[
nonumberlist, %keine Seitenzahlen anzeigen
acronym,      %ein Abkürzungsverzeichnis erstellen
toc]          %Einträge im Inhaltsverzeichnis      
{glossaries}
\usepackage{cite}
%Glossar einschalten
\makeglossaries


%fertigen Kopfzeilenstil aktivieren
\makroFH-Kopfzeilenstil

%Zeilenabstand * 1.25 (default)
\renewcommand{\baselinestretch}{1.25}\normalsize
%(Kommentar entfernen, um Zeilenabstand
% auf 1,5-fache Groesse zu ueberschreiben)
%\renewcommand{\baselinestretch}{1.50}\normalsize

%Startpunkt jedes LaTex-Dokuments
\begin{document}

%<--- Benutze Glossar
%Befehle für Symbole
\newglossaryentry{symb:Pi}{
name=$\pi$,
description={Die Kreiszahl.},
}
\newglossaryentry{symb:Phi}{
name=$\varphi$,
description={Ein beliebiger Winkel.}
}
\newglossaryentry{symb:Lambda}{
name=$\lambda$,
description={Eine beliebige Zahl, mit der der nachfolgende Ausdruck
multipliziert wird.}
}

%Akronyme
\newacronym{MS}{MS}{Microsoft}
\newacronym{CD}{CD}{Compact Disc}
%Eine Akronym mit Glossareintrag
\newacronym{AD}{AD}{Active Directory\protect\glsadd{glos:AD}}

%Glossareinträge
\newglossaryentry{glos:AD}{
name=Active Directory,
description={Active Directory ist in einem Windows 2000/" "Windows
Server 2003-Netzwerk der Verzeichnisdienst, der die zentrale
Organisation und Verwaltung aller Netzwerkressourcen erlaubt. Es
ermöglicht den Benutzern über eine einzige zentrale Anmeldung den
Zugriff auf alle Ressourcen und den Administratoren die zentral
organisierte Verwaltung, transparent von der Netzwerktopologie und
den eingesetzten Netzwerkprotokollen. Das dafür benötigte
Betriebssystem ist entweder Windows 2000 Server oder
Windows Server 2003, welches auf dem zentralen
Domänencontroller installiert wird. Dieser hält alle Daten des
Active Directory vor, wie z.B. Benutzernamen und
Kennwörter.}
}
\newglossaryentry{glos:AntwD}{name=Antwortdatei, description={Informationen zum
Installieren einer Anwendung oder des Betriebssystems.}}



\glsaddall 
%Benutze Glossar --->

%Verwende Muster-Titelseite 
%Titelseite
% Seitennummer aus
\thispagestyle{empty}
\begin{titlepage}
	\vspace{3cm}

%%\begin{center}
%%	\Huge	
%%	HOCHSCHULE LANDSHUT \\
%%	\Large
%%	FAKULTÄT INFORMATIK
%%\end{center}

%%\vspace{1cm}

\begin{center}
	\includegraphics[scale=0.8]{Titelseite/hl-logo.pdf}  
\end{center}

\vspace{2.5cm}

\begin{center}
  \Large FAKULTÄT INFORMATIK
\end{center}

\vspace{1cm}
\begin{center}
	\Huge
	\textbf{Bachelorarbeit/Masterarbeit}\\
\end{center}

\vspace{1cm}

\begin{center}
	\Large
	\textsc{Dies ist der Titel meiner Bacherlor/Masterarbeit}\\
\end{center}

\vspace{1.5cm}

\begin{center}
	\Large
	Vorname Nachname
\end{center}

\vspace{4cm}
\begin{center}
	\large
	Betreuer: Prof. Max Mustermann
\end{center}

\end{titlepage}

%Füge leere Seite ein (optional)
\clearpage 
\ifodd\count0\else 
\thispagestyle{empty} 
\hbox{}\newpage 
\fi 
%Verwende Muster-Erklärung zur selbstständigen Arbeit 
%Erklärung
\thispagestyle{empty}
\vspace{15mm}
\begin{center}
\textbf{\underline{ERKLÄRUNG ZUR BACHELORARBEIT}}
\end{center}
\vspace{25mm}
\begin{center}
\large
%Name, Vorname des Studierenden:  \\
\large
Nachname, Vorname
\end{center}
\vspace{25mm}

\begin{center}
\huge
Hochschule Landshut \\
Fakultät Informatik 
\end{center}
\vspace{10mm}

\begin{center}
\large
Hiermit erkläre ich, dass ich die Arbeit selbständig 
verfasst, noch nicht anderweitig für Prüfungszwecke 
vorgelegt, keine anderen als die angegebenen Quellen 
oder Hilfsmittel benützt, sowie wörtliche und sinngemäße 
Zitate als solche gekennzeichnet habe.  \\
\end{center}
\vspace{55mm}

\begin{center}
....................\hspace{40mm}....................................................\\

(Datum)\hspace{47mm}(Unterschrift des Studierenden)
\end{center}

%Füge leere Seite ein (optional)
\clearpage 
\ifodd\count0\else 
\thispagestyle{empty} 
\hbox{}\newpage 
\fi 
%Verwende Muster-Abstract 
\begin{abstract}
\begin{center}
\Huge
\emph{\textbf{Zusammenfassung/Abstract}}
\end{center}
\normalsize
\vspace{15mm}
\textit{Eine Inhaltsangabe oder Zusammenfassung ist eine Übersicht über den wesentlichen Inhalt eines Textes, Filmes oder Ereignisses. Gebräuchliche Formen von Inhaltsangaben sind das Inhaltsverzeichnis, das Abstract und andere Formen dokumentarischer Referate. Auch die englische Bezeichnung Summary ist in wissenschaftlichen Arbeiten üblich. Die Inhaltsangabe ist in der DIN-Norm DIN 1426 genormt.}
\end{abstract}


%Liste Menüpunkte als Inhaltsverzeichnis
\tableofcontents
\setcounter{page}{1}

%Jedes Kapitel in eigener Ordnerstruktur 
%1.Kapitel ist die Zusammenfassung
\input{Kapitel2/Kapitel2-Text/Kapitel2-Text}
\input{Kapitel3/Kapitel3-Text/Kapitel3-Text}
\input{Kapitel4/Kapitel4-Text/Kapitel4-Text}
\input{Kapitel5/Kapitel5-Text}

% <--- BibTex mit Stil alpha
\bibliography{Bibliothek}{}
\bibliographystyle{alpha}
% BibTex mit Stil alpha --->
 
% Abbildungsverzeichnis, Tabellenverzeichnis und Glossar ausgeben
\listoffigures
\listoftables
\printglossaries 

%Ende jedes LaTex-Dokuments
\end{document}

